% Reading KL Divergence as a Subtraction -- textbook-style chapter.
% Companion to the blog post of the same name (October 2026).
% Build: tectonic kl-divergence-as-a-subtraction.tex   (or latexmk -pdf)
% Every figure is TikZ/pgfplots in this file; evidence_paths.py generates the
% sample paths in Figure A.1.
\documentclass[11pt,oneside]{book}

\usepackage[letterpaper,top=1.05in,bottom=1.1in,inner=1.15in,outer=1.15in,headheight=14pt]{geometry}
\usepackage{amsmath,amsthm,mathtools}
\usepackage{newpxtext,newpxmath}
\usepackage[dvipsnames]{xcolor}
\usepackage{microtype}
\usepackage{booktabs,tabularx,array}
\usepackage{enumitem}
\usepackage{tikz}
\usepackage{pgfplots}
\usepackage[most]{tcolorbox}
\usepackage[font=small,labelfont={bf,sf},format=hang,margin=6pt]{caption}
\usepackage{subcaption}
\usepackage{fancyhdr}
\usepackage{titlesec}
\usepackage[colorlinks=true,linkcolor=pcolor!80!black,citecolor=pcolor!80!black,urlcolor=pcolor!80!black]{hyperref}

\usetikzlibrary{arrows.meta,decorations.pathreplacing,calc,positioning,shapes.geometric}
\pgfplotsset{compat=1.18}
\usepgfplotslibrary{fillbetween}

% ---------- colours: p (target) blue, q (baseline) orange, displacement purple
\definecolor{pcolor}{RGB}{31,78,140}
\definecolor{qcolor}{RGB}{196,96,10}
\definecolor{dcolor}{RGB}{122,46,142}
\definecolor{inkgray}{RGB}{110,110,110}
\definecolor{pos}{RGB}{58,140,90}
\definecolor{neg}{RGB}{190,60,60}

% ---------- notation
\newcommand{\kl}[1]{\mathbin{-^{\mathrm{KL}}_{#1}}}       % p \kl{p} q
\newcommand{\DKL}{D_{\mathrm{KL}}}
\newcommand{\E}{\mathbb{E}}
\newcommand{\pp}{{\color{pcolor}p}}
\newcommand{\qq}{{\color{qcolor}q}}
\newcommand{\X}{\mathcal{X}}

% ---------- theorem-like environments
\theoremstyle{definition}
\newtheorem{definition}{Definition}[chapter]
\newtheorem{example}[definition]{Example}
\newtheorem{exercise}{Exercise}[chapter]
\theoremstyle{plain}
\newtheorem{proposition}[definition]{Proposition}
\theoremstyle{remark}
\newtheorem{remark}[definition]{Remark}

\tcbset{
  keybox/.style={enhanced,breakable,colback=dcolor!4,colframe=dcolor!70,boxrule=0.6pt,
    arc=2pt,left=8pt,right=8pt,top=6pt,bottom=6pt,fonttitle=\sffamily\bfseries\small,
    coltitle=white,colbacktitle=dcolor!75},
  notebox/.style={enhanced,breakable,colback=pcolor!3,colframe=pcolor!45,boxrule=0.5pt,
    arc=2pt,left=8pt,right=8pt,top=5pt,bottom=5pt,fonttitle=\sffamily\bfseries\small,
    coltitle=pcolor!90!black,colbacktitle=pcolor!10}
}

% ---------- tikz styles
\tikzset{
  >={Stealth[length=2.6mm,width=1.8mm]},
  pt/.style={circle,fill=#1,inner sep=1.5pt},
  vec/.style={->,thick,inkgray},
  disp/.style={->,very thick,dcolor},
  back/.style={->,thick,dashed,inkgray},
  lbl/.style={font=\small},
}

% ---------- headings
\titleformat{\chapter}[display]
  {\normalfont\sffamily}{\large\scshape\color{dcolor}Chapter \thechapter}{4pt}
  {\LARGE\bfseries}[\vspace{2pt}{\color{dcolor}\titlerule[0.8pt]}]
\titlespacing*{\chapter}{0pt}{0pt}{18pt}
\titleformat{\section}{\normalfont\sffamily\large\bfseries}{\color{dcolor}\thesection}{0.8em}{}
\titleformat{\subsection}{\normalfont\sffamily\bfseries}{\color{dcolor}\thesubsection}{0.8em}{}

\pagestyle{fancy}
\fancyhf{}
\fancyhead[L]{\small\sffamily\color{inkgray}\nouppercase{\leftmark}}
\fancyhead[R]{\small\sffamily\color{inkgray}Brando Miranda}
\fancyfoot[C]{\small\sffamily\thepage}
\renewcommand{\headrulewidth}{0.3pt}
\renewcommand{\chaptermark}[1]{\markboth{\thechapter.\ #1}{}}
\fancypagestyle{plain}{\fancyhf{}\fancyfoot[C]{\small\sffamily\thepage}\renewcommand{\headrulewidth}{0pt}}

\setlist{itemsep=2pt,topsep=4pt}
\setlength{\parskip}{2pt plus 1pt}

\begin{document}

\chapter{Reading KL Divergence as a Subtraction}

\begin{tcolorbox}[notebox]
\small\textit{Brando Miranda, October 2026.} A textbook-style companion to the blog post
\href{https://cs.stanford.edu/people/brando9/2026/10/17/kl-divergence-as-a-subtraction.html}{\emph{Reading KL Divergence as a Subtraction}}.
The figures are redrawn from my handwritten notes.

\medskip
\textbf{In this chapter.} We build one grammar for the Kullback--Leibler divergence. Read
$\DKL(p\Vert q)$ as $p \kl{p} q$: the divergence of $p$ \emph{from} $q$, measured on samples from $p$.
Unrolled, it is literally a subtraction, $\E_{X\sim p}[\log p(X)-\log q(X)]$. The pointwise gap is
antisymmetric like a vector, the average is not, and \emph{all} of KL's asymmetry comes from which
distribution we average under.
\end{tcolorbox}

\section{Why KL feels arbitrary}

I have been doing machine learning since 2012, and KL divergence has shown up in my losses, my bounds,
my RL regularizers and my variational objectives. Every time it wore a different costume:
\begin{itemize}
  \item the \textbf{excess bits} paid for coding data from $p$ with a code built for $q$;
  \item \textbf{cross-entropy minus entropy}, $H_p(q)-H_p(p)$, with minus signs to babysit;
  \item a \textbf{divergence}, which is ``like a distance'' until it isn't.
\end{itemize}
Each costume is correct. The trouble is having to switch mental models with each textbook. The goal of
this chapter is to replace them with \emph{one notation and one picture} from which every classical formula
falls out. In an earlier post on score matching I subscripted the Fisher divergence with the distribution
that carries the expectation, and promised to do the same for KL. Here that subscript turns out to carry the
whole story of the asymmetry.

\section{The displacement we already trust}

Take two points $A$ and $B$. The vector $B-A$ is the \textbf{displacement of $B$ from $A$}: start at $A$,
end at $B$. Whatever we subtract is the baseline, the place we start from. Its length is the ordinary
distance $d(A,B)=\lVert B-A\rVert_2$. Figure~\ref{fig:displacement}(a) draws it the way I first sketched it:
subtracting $A$ means walking back along $-A$, then out along $B$.

\begin{figure}[t]
\centering
\begin{subfigure}[b]{0.47\textwidth}\centering
\begin{tikzpicture}[scale=1.05]
  \coordinate (O) at (0,0); \coordinate (A) at (3.4,-0.15); \coordinate (B) at (3.6,2.1);
  \draw[vec] (O) -- (B) node[pos=0.55,above left,lbl] {$B$};
  \draw[vec] (O) -- (A) node[pos=0.55,above,lbl] {$A$};
  \draw[back] ($(A)+(0,-0.42)$) -- ($(O)+(0,-0.42)$) node[pos=0.5,below,lbl] {$-A$ \ (undo the baseline)};
  \draw[disp] (A) -- (B) node[midway,right,lbl,text=dcolor] {$B-A$};
  \node[pt=black] at (O) {}; \node[below left,lbl] at (O) {$0$};
  \node[pt=inkgray] at (A) {}; \node[right,lbl] at (A) {$A$ (start)};
  \node[pt=dcolor] at (B) {};  \node[right,lbl] at (B) {$B$ (end)};
\end{tikzpicture}
\caption{$B-A$: the displacement of $B$ from $A$.}
\end{subfigure}\hfill
\begin{subfigure}[b]{0.47\textwidth}\centering
\begin{tikzpicture}[scale=1.05]
  \coordinate (O) at (0,0); \coordinate (Q) at (3.4,-0.15); \coordinate (P) at (3.6,2.1);
  \draw[vec,pcolor!70] (O) -- (P);
  \draw[vec,qcolor!80] (O) -- (Q);
  \draw[back] ($(Q)+(0,-0.42)$) -- ($(O)+(0,-0.42)$) node[pos=0.5,below,lbl] {``$-q$'',\ i.e.\ $\kl{p} q$};
  \draw[disp] (Q) -- (P) node[midway,right,lbl,text=dcolor] {$p \kl{p} q$};
  \node[pt=black] at (O) {};
  \node[pt=qcolor] at (Q) {}; \node[right,lbl] at (Q) {$\qq$ (baseline)};
  \node[pt=pcolor] at (P) {}; \node[right,lbl] at (P) {$\pp$ (target)};
\end{tikzpicture}
\caption{The same grammar for distributions: $q$ ``goes to'' $p$.}
\end{subfigure}
\caption{Displacement grammar, redrawn from page~1 of my notes. Panel (b) is the reading I propose for
KL: $q \xrightarrow[p]{\mathrm{KL}} p$, start at the baseline $q$ and arrive at the target $p$. On my first pass I
had drawn the arrow the other way, $p\to q$; with the $B-A$ grammar it points from $q$ to $p$.}
\label{fig:displacement}
\end{figure}

Everything we like about distances is visible in this picture. The length is nonnegative and zero only when
$A=B$; it is symmetric; it obeys the triangle inequality. Keep one fact in mind for later:
\begin{tcolorbox}[keybox]
\centering The \textbf{displacement is antisymmetric}, $A-B=-(B-A)$; the \textbf{length is symmetric}, because
the norm forgets the sign.
\end{tcolorbox}
That split is what explains KL.

\section{Distributions as points, and the notation}

Let $p$ and $q$ be distributions on the same space $\X$. On a finite space, $p\in[0,1]^{\X}$ is a vector
with one coordinate per outcome $x$, its value $p(x)$, normalized so that $\sum_x p(x)=1$. In the continuous
case it is an ``infinite vector'' indexed by $x$, a density with $\int_{\X} p(x)\,dx=1$. This is an analogy
rather than a theorem, but it makes the right question obvious: \emph{what is the displacement of $p$ from $q$?}

\begin{definition}[Displacement notation for KL]\label{def:notation}
For distributions $p,q$ on $\X$ with $p\ll q$,
\[
\DKL(p\Vert q)\;=\;D^{\mathrm{KL}}_p(p\Vert q)\;=:\;p\kl{p}q
\;=\;\Delta^{\mathrm{KL}}_p(p,q)\;=\;\Delta^{\mathrm{KL}}_p(q\to p)\;=\;\Delta^{\mathrm{KL}}_p(p\leftarrow q).
\]
Read it as \textbf{the divergence of $p$ from $q$, measured on samples from $p$}. The superscript names the
kind of subtraction; the subscript names the measure we average under, the ``ruler''. The baseline is $q$
and the destination is $p$, exactly as $B-A$ has baseline $A$. Usually $p$ is the target (the data) and $q$
the model that approximates it.
\end{definition}

``Divergence'' here just means an unsymmetric distance, difference or displacement. Both arrows in
Definition~\ref{def:notation} say the same thing read in opposite directions: start at $q$, arrive at $p$.
Figure~\ref{fig:page6} is the small spiral-pad version of the same idea, where the reading in bits sits next
to the triangle.

\begin{figure}[t]
\centering
\begin{subfigure}[b]{0.44\textwidth}\centering
\begin{tikzpicture}
  \draw[inkgray!60] (-0.3,0) -- (4.6,0);
  \node[pt=inkgray] (b) at (0.6,0) {}; \node[below=2pt,lbl] at (b) {$B$};
  \node[pt=dcolor]  (a) at (3.8,0) {}; \node[below=2pt,lbl] at (a) {$A$};
  \draw[disp] (0.6,0.35) -- (3.8,0.35) node[midway,above,lbl,text=dcolor] {$A-B$};
  \node[lbl,align=center] at (2.2,-1.0) {distance of $A$ from $B$};
  \node at (2.2,-1.75) {};
\end{tikzpicture}
\caption{On a line: $A-B$ is how far $A$ sits from $B$.}
\end{subfigure}\hfill
\begin{subfigure}[b]{0.52\textwidth}\centering
\begin{tikzpicture}
  \coordinate (O) at (0,0.2); \coordinate (Q) at (2.6,0); \coordinate (P) at (2.75,2.0);
  \draw[inkgray] (O) -- (Q); \draw[inkgray] (O) -- (P);
  \draw[disp] (Q) -- (P);
  \node[pt=qcolor] at (Q) {}; \node[below,lbl] at (Q) {$\qq$};
  \node[pt=pcolor] at (P) {}; \node[above,lbl] at (P) {$\pp$};
  \draw[decorate,decoration={brace,amplitude=6pt},thick] ($(P)+(0.35,0)$) -- ($(Q)+(0.35,0)$)
     node[midway,right=8pt,lbl,align=left] {excess bits to represent\\ $p$ using $q$, when\\ sampling from $p$};
\end{tikzpicture}
\caption{$p\kl{p}q$, the ``divergence'' of $p$ from $q$ w.r.t.\ $p$.}
\end{subfigure}
\caption{Redrawn from page~6 of my notes. The one-dimensional picture (a) is the one everyone trusts. In (b)
the vertical edge is the KL divergence and the brace gives its classical reading:
$p^{*}\kl{p^{*}}\hat q=H_{p^{*}}(\hat q)-H_{p^{*}}(p^{*})$.}
\label{fig:page6}
\end{figure}

Put the two pictures side by side and the analogy reads
\begin{equation}\label{eq:analogy}
p\kl{p}q\;=\;\E_{X\sim p}\big[\log p(X)-\log q(X)\big]\;\approx\;\text{``}A-B\text{''},
\end{equation}
the divergence of $p$ from $q$ next to the distance of $A$ from $B$ (Figure~\ref{fig:vectorview}).

\begin{remark}[A slight abuse of notation]
The quotes around ``$A-B$'' matter. The vector $\vec{A}-\vec{B}$ is a displacement; the \emph{distance} is the
scalar $\lVert A-B\rVert_2$, and here ``$A-B$'' is shorthand for it. Likewise $p\kl{p}q$ is not a vector: it is a
scalar, like the Euclidean distance. Unlike that distance it is asymmetric, and the direction lives in the
subscript.
\end{remark}

\begin{figure}[t]
\centering
\begin{tikzpicture}
  % left: distributions
  \begin{scope}
    \coordinate (O) at (0,0.15); \coordinate (Q) at (2.7,0); \coordinate (P) at (2.85,2.3);
    \draw[inkgray] (O) -- (P);
    \draw[back] (Q) -- (O) node[midway,below,lbl] {$\kl{p}$};
    \draw[disp] (Q) -- (P) node[midway,right,lbl,text=dcolor] {$p\kl{p}q$};
    \node[pt=qcolor] at (Q) {}; \node[below right,lbl] at (Q) {$\qq$};
    \node[pt=pcolor] at (P) {}; \node[above,lbl] at (P) {$\pp$};
    \node[lbl,align=center] at (1.5,-1.0) {divergence (u-dist) of $p$ from $q$\\[1pt] $=\Delta^{\mathrm{KL}}_p(q\to p)$};
  \end{scope}
  \node[align=center] at (5.0,1.1) {\small analogy\\[2pt] {\Large$\approx$}};
  % right: points
  \begin{scope}[xshift=6.6cm]
    \coordinate (O) at (0,0.15); \coordinate (B) at (2.7,0); \coordinate (A) at (2.85,2.3);
    \draw[inkgray] (O) -- (A);
    \draw[back] (B) -- (O) node[midway,below,lbl] {$-B$};
    \draw[disp] (B) -- (A) node[midway,right,lbl,text=dcolor,align=left] {$\vec{A}-\vec{B}$\\ $=\Delta(A-B)$};
    \node[pt=inkgray] at (B) {}; \node[below right,lbl] at (B) {$B$};
    \node[pt=dcolor] at (A) {}; \node[above,lbl] at (A) {$A$};
    \node[lbl,align=center] at (1.5,-1.0) {distance of $A$ from $B$\\[1pt] ``$A-B$'' $\stackrel{?}{=}\lVert A-B\rVert_2$};
  \end{scope}
  \node[draw=dcolor!60,rounded corners=2pt,fill=dcolor!4,font=\footnotesize,align=left,text width=11.2cm,inner sep=5pt]
     at (4.95,-2.45) {\textbf{Note:} a little abuse of notation. $p\kl{p}q$ is not a vector, it is a scalar like
     \mbox{$\lVert A-B\rVert_2$}. The Euclidean distance of $A$ from $B$ is symmetric; the divergence of $p$ from $q$ is not.};
\end{tikzpicture}
\caption{The vector view, redrawn from page~7 of my notes (its final version). Left: the divergence of $p$
from $q$, with the measure $p$ in the subtraction sign. Right: the displacement of $A$ from $B$, with
``$A-B$'' standing for the distance.}
\label{fig:vectorview}
\end{figure}

The same page asks a question in the margin: since we sample from $p$, does $\log p(X)$ come out the bigger
side? \textbf{On average, yes}, by Gibbs' inequality (Proposition~\ref{prop:gibbs}). \textbf{Not sample by sample}: a
particular $x$ drawn from $p$ can still be likelier under $q$ (Figure~\ref{fig:pointwise}).

\section{Unroll it: KL really is a subtraction}

Write out the definition and the subtraction is sitting right there:
\begin{equation}\label{eq:unrolled}
p\kl{p}q\;=\;\E_{X\sim p}\big[\,\log p(X)\;-\;\log q(X)\,\big].
\end{equation}

\begin{definition}[Log-likelihood displacement]
The pointwise log-likelihood displacement of $p$ from $q$ is $\ell(x):=\log p(x)-\log q(x)$.
Then $p\kl{p}q=\E_p[\ell(X)]$.
\end{definition}

Read~\eqref{eq:unrolled} as a procedure (Figure~\ref{fig:loglik}b): sample $x\sim p$; score it under both models;
subtract, which tells you how much likelier $p$ makes this $x$ than $q$ does, in log units (nats with $\ln$, bits
with $\log_2$); average. So $\DKL(p\Vert q)$ is \textbf{the average log-likelihood advantage of $p$ over $q$ on $p$'s
own data}. No codes, and no minus signs to keep track of.

\begin{figure}[t]
\centering
\begin{subfigure}[b]{0.44\textwidth}\centering
\begin{tikzpicture}[scale=1.05]
  \coordinate (O) at (0,0); \coordinate (R) at (3.3,0); \coordinate (T) at (3.3,2.2);
  \draw[vec,pcolor] (O) -- (T) node[midway,above left,lbl,text=pcolor] {$\log p(X)$};
  \draw[vec,qcolor] ($(R)+(0,-0.35)$) -- ($(O)+(0,-0.35)$) node[midway,below,lbl,text=qcolor] {$-\log q(X)$};
  \draw[inkgray!50] (O) -- (R);
  \draw[disp] (R) -- (T) node[midway,right,lbl,text=dcolor,align=left] {$\ell(X)$};
  \node[pt=black] at (O) {}; \node[pt=qcolor] at (R) {}; \node[pt=pcolor] at (T) {};
\end{tikzpicture}
\caption{The triangle in log-likelihood space.}
\end{subfigure}\hfill
\begin{subfigure}[b]{0.52\textwidth}\centering
\begin{tikzpicture}[box/.style={draw=dcolor!70,rounded corners=2pt,fill=dcolor!5,font=\small,
    minimum width=5.6cm,minimum height=0.62cm,align=center},node distance=0.33cm]
  \node[box] (s) {\textbf{1. Sample} \ $x\sim p$};
  \node[box,below=of s] (c) {\textbf{2. Score} \ ${\color{pcolor}\log p(x)}$ and ${\color{qcolor}\log q(x)}$};
  \node[box,below=of c] (d) {\textbf{3. Subtract} \ $\ell(x)=\log p(x)-\log q(x)$};
  \node[box,below=of d] (a) {\textbf{4. Average} \ $\frac1n\sum_i \ell(x_i)\to p\kl{p}q$};
  \foreach \u/\v in {s/c,c/d,d/a} \draw[->,dcolor!80] (\u) -- (\v);
\end{tikzpicture}
\caption{KL as a recipe.}
\end{subfigure}
\caption{Redrawn from page~5 of my notes, the page where it clicked. (a) Walking out along $\log p(X)$ and back
along $-\log q(X)$ leaves the displacement $\ell(X)=\log p(X)-\log q(X)$, just as $B-A=B+(-A)$. (b) Sample, compute
log-probabilities, compare, average.}
\label{fig:loglik}
\end{figure}

\begin{proposition}[Gibbs' inequality]\label{prop:gibbs}
$p\kl{p}q\ge0$, with equality if and only if $p=q$ almost everywhere.
\end{proposition}
\begin{proof}
By Jensen's inequality for the concave $\log$,
\[
-\E_p[\ell(X)]=\E_p\!\left[\log\frac{q(X)}{p(X)}\right]\le\log\E_p\!\left[\frac{q(X)}{p(X)}\right]
=\log\!\int_{\{p>0\}} q(x)\,dx\le\log1=0,
\]
with equality only when $q/p$ is constant on the support of $p$, that is, when $p=q$.
\end{proof}

For a particular $x$, $\ell(x)$ can be negative: $q$ may like some of $p$'s samples more than $p$ does. On average,
on $p$'s own turf, $p$ always wins or ties. Figure~\ref{fig:pointwise} shows both facts at once.

\begin{figure}[t]
\centering
\begin{tikzpicture}
\begin{axis}[name=top,width=12.4cm,height=4.0cm,xmin=-4.5,xmax=4.5,ymin=0,ymax=0.45,
  axis lines=left,xticklabels={},ytick=\empty,ylabel={density},ylabel style={font=\small},
  legend style={draw=none,font=\small,at={(0.98,0.95)},anchor=north east},clip=false]
  \addplot[pcolor,very thick,domain=-4.5:4.5,samples=160]{exp(-x^2/2)/sqrt(2*pi)};
  \addlegendentry{$p=\mathcal N(0,1)$}
  \addplot[qcolor,very thick,dashed,domain=-4.5:4.5,samples=160]{exp(-x^2/8)/(2*sqrt(2*pi))};
  \addlegendentry{$q=\mathcal N(0,2^2)$}
\end{axis}
\begin{axis}[at={(top.south west)},anchor=north west,yshift=-0.35cm,width=12.4cm,height=5.2cm,
  xmin=-4.5,xmax=4.5,ymin=-6.0,ymax=2.0,axis lines=left,xlabel={$x$},
  ytick={-6,-4,-2,0},ylabel={nats},ylabel style={font=\small},clip=true,
  legend style={draw=none,font=\small,at={(0.02,0.04)},anchor=south west,fill=white,fill opacity=0.9,text opacity=1}]
  \addplot[name path=ell,dcolor,very thick,domain=-4.5:4.5,samples=160]{ln(2)-3*x^2/8};
  \addlegendentry{$\ell(x)=\log p(x)-\log q(x)$}
  \addplot[black,thick,densely dotted,domain=-4.5:4.5,samples=160]{exp(-x^2/2)/sqrt(2*pi)*(ln(2)-3*x^2/8)};
  \addlegendentry{integrand $p(x)\,\ell(x)$}
  \path[name path=zero] (axis cs:-4.5,0) -- (axis cs:4.5,0);
  \addplot[pos!25,forget plot] fill between[of=ell and zero,soft clip={domain=-1.36:1.36}];
  \addplot[neg!18,forget plot] fill between[of=ell and zero,soft clip={domain=1.36:4.5}];
  \addplot[neg!18,forget plot] fill between[of=ell and zero,soft clip={domain=-4.5:-1.36}];
  \draw[inkgray] (axis cs:-4.5,0) -- (axis cs:4.5,0);
  \node[font=\small,text=pos!70!black,anchor=south] at (axis cs:0,0.9) {$p$ wins here: $\ell>0$};
  \node[font=\small,text=neg!80!black,align=center] at (axis cs:3.0,-1.3) {tails: $\ell<0$,\\ $q$ likes these $x$ more};
\end{axis}
\end{tikzpicture}
\caption{Pointwise versus average, for $p=\mathcal N(0,1)$ and $q=\mathcal N(0,2^2)$, where
$\ell(x)=\ln2-\tfrac38x^2$. Samples from $p$ with $|x|>\sqrt{(8\ln2)/3}\approx1.36$ (about 17\% of them) are
\emph{likelier} under $q$. The dotted integrand still has positive total area,
$p\kl{p}q=\ln2-\tfrac38\approx0.32$ nats, as Gibbs' inequality promises.}
\label{fig:pointwise}
\end{figure}

\begin{example}[Sometimes it literally is a distance]\label{ex:equalvar}
Let $p=\mathcal N(\mu_p,\sigma^2)$ and $q=\mathcal N(\mu_q,\sigma^2)$. Then
$\ell(x)=\frac{(x-\mu_q)^2-(x-\mu_p)^2}{2\sigma^2}$, and averaging under $p$ gives
\[
p\kl{p}q=\frac{(\mu_p-\mu_q)^2}{2\sigma^2},
\]
half the squared Euclidean distance between the means, in units of $\sigma$ (Figure~\ref{fig:equalvar}). In the
simplest case the displacement picture is not a mnemonic; it is exactly the old distance.
\end{example}

\begin{figure}[t]
\centering
\begin{tikzpicture}
\begin{axis}[width=10.5cm,height=4.3cm,xmin=-3.4,xmax=5.4,ymin=0,ymax=0.56,axis lines=left,
  ytick=\empty,xtick={0,2},xticklabels={$\mu_q$,$\mu_p$},clip=false]
  \addplot[qcolor,very thick,dashed,domain=-3.4:5.4,samples=140]{exp(-x^2/2)/sqrt(2*pi)};
  \addplot[pcolor,very thick,domain=-3.4:5.4,samples=140]{exp(-(x-2)^2/2)/sqrt(2*pi)};
  \draw[inkgray,densely dotted] (axis cs:0,0) -- (axis cs:0,0.4);
  \draw[inkgray,densely dotted] (axis cs:2,0) -- (axis cs:2,0.4);
  \draw[disp] (axis cs:0,0.45) -- (axis cs:2,0.45);
  \node[font=\small,text=dcolor,anchor=south] at (axis cs:1,0.46) {$p\kl{p}q=\frac{(\mu_p-\mu_q)^2}{2\sigma^2}$};
  \node[font=\small,text=qcolor] at (axis cs:-1.9,0.27) {$q$};
  \node[font=\small,text=pcolor] at (axis cs:3.9,0.27) {$p$};
\end{axis}
\end{tikzpicture}
\caption{Example~\ref{ex:equalvar}: with equal variances the KL divergence is the displacement of the means
from baseline $\mu_q$ to target $\mu_p$, squared and halved ($\sigma=1$ and $\mu_p-\mu_q=2$ give $2$ nats).
Here it is symmetric; the asymmetry appears as soon as the shapes differ.}
\label{fig:equalvar}
\end{figure}

\section{Where the asymmetry lives}

Swap $p$ and $q$ in the pointwise displacement and it flips sign, exactly like a vector:
$\log q(x)-\log p(x)=-\ell(x)$. So the reverse divergence is
\begin{equation}\label{eq:reverse}
q\kl{q}p\;=\;\E_{X\sim q}\big[-\ell(X)\big].
\end{equation}
\begin{tcolorbox}[keybox,title=Where the asymmetry comes from]
Same pointwise displacement, opposite sign, \textbf{different averaging distribution}. In Euclidean space the norm
forgets the sign and symmetry comes for free. In KL the ``norm'' is an average taken under one of the two
endpoints; change the endpoint you sample from and you change the number. The subscript in $\kl{p}$ records
exactly this.
\end{tcolorbox}

\begin{example}[Narrow versus wide]
Let $p=\mathcal N(0,1)$ and $q=\mathcal N(0,2^2)$. Then
\[
p\kl{p}q=\ln2-\tfrac38\approx0.32\text{ nats},\qquad q\kl{q}p=\tfrac32-\ln2\approx0.81\text{ nats}.
\]
On $p$'s data the wide $q$ is a little diffuse but never shocked, so the cost is low. On $q$'s data the narrow $p$ keeps
meeting samples in the tails that it thought nearly impossible, so the cost is high.
\end{example}

This is the fact behind ``forward KL is mode-covering, reverse KL is mode-seeking'', and behind the different
behaviour of maximum likelihood (forward) and variational inference (reverse). Fit a single Gaussian $q$ to a
two-mode $p$ in each direction and the subscript decides everything (Figure~\ref{fig:fwdrev}).

\begin{figure}[t]
\centering
\pgfmathdeclarefunction{gauss}{2}{\pgfmathparse{exp(-((x-#1)^2)/(2*#2^2))/(#2*sqrt(2*pi))}}
\begin{tikzpicture}
\begin{axis}[name=L,width=7.0cm,height=4.8cm,xmin=-5,xmax=5,ymin=0,ymax=0.75,axis lines=left,
  ytick=\empty,xtick={-4,-2,0,2,4},title={\small Forward: $\min_q\, p\kl{p}q$\\\footnotesize averaged under $p$},
  title style={align=center,yshift=-3pt}]
  \addplot[draw=none,fill=pcolor!22,domain=-5:5,samples=200]{0.35*gauss(-2,0.6)+0.65*gauss(2,0.6)}\closedcycle;
  \addplot[pcolor!80,thick,domain=-5:5,samples=200]{0.35*gauss(-2,0.6)+0.65*gauss(2,0.6)};
  \addplot[qcolor,very thick,domain=-5:5,samples=160]{gauss(0.6,2.0)};
  \node[font=\small,text=qcolor,fill=white,fill opacity=0.8,text opacity=1,inner sep=1pt] at (axis cs:-1.2,0.5) {covers both modes};
\end{axis}
\begin{axis}[at={(L.east)},anchor=west,xshift=0.9cm,width=7.0cm,height=4.8cm,xmin=-5,xmax=5,ymin=0,ymax=0.75,
  axis lines=left,ytick=\empty,xtick={-4,-2,0,2,4},title={\small Reverse: $\min_q\, q\kl{q}p$\\\footnotesize averaged under $q$},
  title style={align=center,yshift=-3pt}]
  \addplot[draw=none,fill=pcolor!22,domain=-5:5,samples=200]{0.35*gauss(-2,0.6)+0.65*gauss(2,0.6)}\closedcycle;
  \addplot[pcolor!80,thick,domain=-5:5,samples=200]{0.35*gauss(-2,0.6)+0.65*gauss(2,0.6)};
  \addplot[qcolor,very thick,domain=-5:5,samples=200]{gauss(2.0,0.6)};
  \draw[->,inkgray] (axis cs:-0.4,0.6) -- (axis cs:1.45,0.6);
  \node[font=\small,text=qcolor,align=right,anchor=east] at (axis cs:-0.4,0.6) {locks onto\\ one mode};
  \node[font=\small,text=inkgray,align=center] at (axis cs:-2,0.34) {ignored};
\end{axis}
\end{tikzpicture}
\caption{The best single Gaussian $q$ (orange) for the two-mode $p=0.35\,\mathcal N(-2,0.6^2)+0.65\,\mathcal N(2,0.6^2)$
(shaded). Left: averaged under $p$, any place where $p$ has mass and $q$ does not costs a huge $\ell(x)$, so $q$
stretches to cover everything; the optimum matches moments, $q=\mathcal N(0.6,2.0^2)$, at $0.56$ nats. Right:
averaged under $q$, $q$ is scored only where it puts its own samples, so it may ignore a whole mode; the optimum
is $q\approx\mathcal N(2.0,0.6^2)$, at $0.43\approx-\ln0.65$ nats.}
\label{fig:fwdrev}
\end{figure}

Koller and Friedman's treatment of M- and I-projections \cite[Sec.~8.5]{koller2009} is the cleanest account of
this I know. If a symmetric quantity is needed, add the two directions (Jeffreys' divergence):
\begin{equation}\label{eq:jeffreys}
p\kl{p}q\;+\;q\kl{q}p\;=\;\int\big(p(x)-q(x)\big)\big(\log p(x)-\log q(x)\big)\,dx .
\end{equation}
A difference in probability space times a difference in log space: both subtractions in one integrand.

\section{How KL is, and is not, like a distance}

Table~\ref{tab:distance} sets KL beside the quantity we started from.

\begin{table}[h]
\centering\small
\renewcommand{\arraystretch}{1.25}
\begin{tabularx}{\textwidth}{@{}>{\raggedright\arraybackslash}p{3.6cm}>{\raggedright\arraybackslash}X>{\raggedright\arraybackslash}X@{}}
\toprule
 & Half squared Euclidean, $\tfrac12\lVert B-A\rVert^2$ & KL, $p\kl{p}q$\\
\midrule
Pointwise displacement & $B-A$ & $\log p(x)-\log q(x)$\\
Swap the endpoints & flips the displacement's sign & flips the displacement's sign\\
Nonnegative, zero only at equality & yes & yes (Gibbs)\\
Symmetric & yes & no: the average is under one endpoint\\
Triangle inequality & no (the unsquared distance has it) & no\\
Bregman divergence of & $\tfrac12\lVert x\rVert^2$ & negative entropy $\sum_x p(x)\log p(x)$\\
Locally, near equality & itself & $\tfrac12\,\delta^{\top}F(\theta)\,\delta$, $F$ the Fisher information\\
\bottomrule
\end{tabularx}
\caption{KL next to half the squared Euclidean distance.}
\label{tab:distance}
\end{table}

The honest summary: \textbf{KL is not a distance, but it behaves like half a squared distance.} Both belong to the
family of Bregman divergences,
\[
D_F(p,q)\;=\;F(p)-F(q)-\langle\nabla F(q),\,p-q\rangle,
\]
the gap at $p$ between a convex function $F$ and its tangent drawn at the baseline $q$ (Figure~\ref{fig:bregman}).
With $F(x)=\tfrac12\lVert x\rVert^2$ this is $\tfrac12\lVert p-q\rVert^2$; with the negative entropy it is
$\DKL(p\Vert q)$. Nearby distributions see a quadratic form,
$\DKL(p_\theta\Vert p_{\theta+\delta})\approx\tfrac12\delta^{\top}F(\theta)\delta$, so locally KL is a squared length in the
Fisher metric. It even has a Pythagorean identity for projections onto nice families \cite{csiszar1975}. What it gives up is
symmetry, and we now know exactly why.

\begin{figure}[t]
\centering
\begin{tikzpicture}
\begin{axis}[width=9.6cm,height=5.6cm,xmin=0,xmax=4.2,ymin=-0.6,ymax=4.4,axis lines=left,
  xtick={1.1,3.4},xticklabels={$q$ (baseline),$p$ (target)},ytick=\empty,clip=false]
  \addplot[black,very thick,domain=0.05:4.1,samples=120]{0.35*x^2-0.2*x+0.3};
  \node[font=\small] at (axis cs:3.95,4.0) {$F$};
  % tangent at q=1.1: F(1.1)=0.5035, F'(1.1)=0.57
  \addplot[inkgray,thick,dashed,domain=0.1:4.0]{0.5035+0.57*(x-1.1)};
  \node[font=\small,text=inkgray,anchor=west] at (axis cs:4.0,2.17) {tangent at $q$};
  \addplot[only marks,mark=*,mark size=1.8pt,qcolor] coordinates {(1.1,0.5035)};
  \addplot[only marks,mark=*,mark size=1.8pt,pcolor] coordinates {(3.4,3.666)};
  \addplot[only marks,mark=*,mark size=1.6pt,inkgray] coordinates {(3.4,1.8145)};
  \draw[inkgray,densely dotted] (axis cs:1.1,-0.6) -- (axis cs:1.1,0.5035);
  \draw[inkgray,densely dotted] (axis cs:3.4,-0.6) -- (axis cs:3.4,1.8145);
  \draw[disp,<->] (axis cs:3.4,1.86) -- (axis cs:3.4,3.62);
  \node[font=\small,text=dcolor,anchor=west] at (axis cs:3.45,2.6) {$D_F(p,q)$};
\end{axis}
\end{tikzpicture}
\caption{A Bregman divergence is the gap at the target $p$ between a convex $F$ and its tangent at the
baseline $q$. Moving the tangent point changes the gap, which is why such divergences are asymmetric in general.
Half the squared Euclidean distance ($F=\tfrac12\lVert x\rVert^2$) and KL ($F$ = negative entropy) are both of this form.}
\label{fig:bregman}
\end{figure}

\section{Why I would teach the log-likelihood view first}

\paragraph{The signs make sense.} The bits view counts \emph{cost}, so it is built from $-\log$ terms, and KL becomes
cost minus cost with sign flips to babysit. The log-likelihood view counts \emph{plausibility}: bigger is better, and
KL is the plausibility gap. It is the same gap on a mirrored axis (Figure~\ref{fig:axes}).

\begin{figure}[t]
\centering
\begin{tikzpicture}[x=1.15cm]
  % top: surprisal axis, cost
  \node[anchor=west,font=\small\sffamily\bfseries] at (-0.2,1.85) {Surprisal $I(\cdot)=-\log(\cdot)$: cost, bigger is worse};
  \draw[->,thick] (0,0.5) -- (10,0.5);
  \foreach \x/\t in {0/$0$,3.5/${I(p)}$,7.0/${I(q)}$} {\draw (\x,0.4) -- (\x,0.6); \node[below,font=\small] at (\x,0.4) {\t};}
  \node[pt=pcolor] at (3.5,0.5) {}; \node[pt=qcolor] at (7.0,0.5) {};
  \draw[decorate,decoration={brace,amplitude=5pt},thick,dcolor] (3.5,0.72) -- (7.0,0.72)
     node[midway,above=5pt,font=\small,text=dcolor] {excess bits $I(q)-I(p)$};
  % bottom: log-likelihood axis, plausibility
  \node[anchor=west,font=\small\sffamily\bfseries] at (-0.2,-0.6) {Log-likelihood $\log(\cdot)$: plausibility, bigger is better};
  \draw[->,thick] (0,-1.9) -- (10,-1.9);
  \foreach \x/\t in {3.0/${\log q}$,6.5/${\log p}$,9.5/$0$} {\draw (\x,-2.0) -- (\x,-1.8); \node[below,font=\small] at (\x,-2.0) {\t};}
  \node[pt=qcolor] at (3.0,-1.9) {}; \node[pt=pcolor] at (6.5,-1.9) {};
  \draw[decorate,decoration={brace,amplitude=5pt},thick,dcolor] (3.0,-1.68) -- (6.5,-1.68)
     node[midway,above=5pt,font=\small,text=dcolor] {gap $\log p-\log q=\ell$};
  
\end{tikzpicture}
\caption{Redrawn from page~4 of my notes, the margin of a book chapter ($p$ the reference distribution we sample
from, $q$ the approximating one). On the surprisal axis KL is $\E_p[I(q)-I(p)]=H_p(q)-H_p(p)$; flip the axis and the
same gap is $\E_p[\log p-\log q]$. Drawn for a typical $x\sim p$, where $p$ wins.}
\label{fig:axes}
\end{figure}

\paragraph{Maximum likelihood falls out in one line.} With $p$ the data distribution and $q_\theta$ the model,
\[
p\kl{p}q_\theta\;=\;\underbrace{\E_p[\log p(X)]}_{\text{independent of }\theta}\;-\;\E_p[\log q_\theta(X)].
\]
Minimizing over $\theta$ is maximizing the expected log-likelihood, and with data we replace $\E_p$ by the
empirical average. \textbf{Maximum likelihood minimizes the forward KL from the data to the model}: we move the
baseline $q_\theta$ toward the data, scored on the data's own samples.

\paragraph{It connects to the rest of statistics.} $\ell(x)$ is the log-likelihood ratio, the central object of
hypothesis testing, Bayesian evidence, importance sampling and policy-gradient RL. Once KL is ``average $\ell$'',
much of ML stops looking like a collection of unrelated tricks (Appendix~\ref{app:ratios}).

\section{Notes and history}

I did not get this view from a course or a textbook. I got it while working through these notes with an AI model:
I asked about the excess-bits formula, and it rewrote KL as $\E_p[\log p-\log q]$ without my asking.

\textbf{It is not new as mathematics.} Kullback and Leibler \cite{kullback1951} introduced the quantity as the mean
information for discrimination between two hypotheses per observation from one of them, which is precisely the
expected log-likelihood ratio, and Kullback \cite{kullback1959} preferred the name \emph{discrimination information}.
Akaike built AIC on the same reading \cite{akaike1973}, and Etz \cite{etz2018} calls KL ``the expected log likelihood
ratio in favor of the true distribution.'' What I could not find in the books on my shelf, those I was assigned or
recommended at UIUC, MIT and Stanford, is any of them \emph{leading} with it as the intuition (Table~\ref{tab:books}).
Nor could I find it folded into one displacement grammar in which the averaging distribution is part of the notation
and explains the asymmetry. That part is my framing.

\begin{table}[h]
\centering\footnotesize
\renewcommand{\arraystretch}{1.2}
\begin{tabularx}{\textwidth}{@{}>{\raggedright\arraybackslash}p{3.0cm}>{\raggedright\arraybackslash}X>{\raggedright\arraybackslash}X>{\raggedright\arraybackslash}p{2.1cm}>{\raggedright\arraybackslash}p{2.3cm}@{}}
\toprule
Book & Expected log-difference or log-ratio? & Sampled log-likelihood form? & Headline intuition & Where\\
\midrule
Hogg, McKean \& Craig \cite{hogg2013} & Yes, unnamed: $\E_{\theta_0}\!\big[\log\frac{f(X;\theta)}{f(X;\theta_0)}\big]$ & Yes: average log-likelihood ratio over the sample & None; KL never named & Thm.~6.1.1 and proof\\
MacKay \cite{mackay2003} & Yes, $\sum_x P(x)\log\frac{P(x)}{Q(x)}$ & Partly: expected weight of evidence, not linked to KL & Wasted bits & Sec.~2.6; ch.~5; Secs.~18.3--18.4\\
Hastie, Tibshirani \& Friedman \cite{hastie2009} & Only as ICA's mutual information & No & Dependence & Sec.~14.7.2\\
Koller \& Friedman \cite{koller2009} & Yes; a proof writes $\E_P[\log P-\log P']$ & Yes & Extra coding cost & Def.~A.5; Prop.~16.1; Sec.~17.2.5\\
Shalev-Shwartz \& Ben-David \cite{ssbd2014} & Yes, inside the log-loss risk & Yes: MLE as ERM with log-loss & Log-loss risk $=$ relative entropy $+$ entropy & Sec.~24.1\\
Goodfellow, Bengio \& Courville \cite{goodfellow2016} & Yes, exactly $\E_{x\sim P}[\log P(x)-\log Q(x)]$ & Yes: MLE as KL to the empirical distribution & Extra bits & Eq.~3.50; eq.~5.60\\
Bishop \& Bishop \cite{bishop2024} & Cross-entropy minus entropy, then $-\int p\ln\frac qp$ & In Ex.~2.34 & Extra information from the wrong code & Sec.~2.5.5\\
Bishop \cite{bishop2006} & Cross-entropy minus entropy, then $-\int p\ln\frac qp$ & Yes: $\frac1N\sum_n\{-\ln q(x_n|\theta)+\ln p(x_n)\}$ & Extra information to transmit & Sec.~1.6.1\\
\bottomrule
\end{tabularx}
\caption{How the books on my shelf introduce KL. The subtraction is in almost all of them; the headline is always bits.}
\label{tab:books}
\end{table}

Why did it take fourteen years? My guess, with no complaint against anyone, is that a class, a talk or a paper has a
syllabus or a page limit, and ``why is it a ratio, and why that sign?'' feels like something one should already know.
So nobody stops to rebuild the foundations. AI changes the economics of asking: a tireless sparring partner that never
makes you feel slow. A curious person who works hard with it can now understand the foundations more deeply than any
course gave them.

\section*{Exercises}
\addcontentsline{toc}{section}{Exercises}

\begin{exercise}
Derive Example~\ref{ex:equalvar}: for $p=\mathcal N(\mu_p,\sigma^2)$ and $q=\mathcal N(\mu_q,\sigma^2)$, show
$p\kl{p}q=(\mu_p-\mu_q)^2/(2\sigma^2)$ by averaging $\ell(x)$ under $p$.
\end{exercise}

\begin{exercise}
For $p=\mathcal N(0,1)$ and $q=\mathcal N(0,s^2)$, compute $p\kl{p}q$ and $q\kl{q}p$ as functions of $s$. For which $s$
are they equal? Check the values $0.32$ and $0.81$ at $s=2$.
\end{exercise}

\begin{exercise}
Prove the Jeffreys identity~\eqref{eq:jeffreys}, and show its integrand is pointwise nonnegative. Why does this not
make each direction nonnegative pointwise?
\end{exercise}

\begin{exercise}
Let $y=g(x)$ with $g$ smooth and invertible, and let $p_Y,q_Y$ be the induced densities. Show that
$\ell$ is unchanged, $\log p_Y(y)-\log q_Y(y)=\log p(x)-\log q(x)$, so $p\kl{p}q$ is the same in both coordinates,
while the pointwise difference $p_Y(y)-q_Y(y)$ is rescaled by the Jacobian.
\end{exercise}

\begin{exercise}
Show that the Gaussian $q$ minimizing $p\kl{p}q$ matches the mean and variance of $p$. Recover the left panel of
Figure~\ref{fig:fwdrev}: $\mu=0.6$, $\sigma=2.0$.
\end{exercise}

\begin{exercise}
In PPO the ratio is $r_t(\theta)=\pi_\theta(a_t\mid s_t)/\pi_{\theta_{\mathrm{old}}}(a_t\mid s_t)$. Which KL do you get by
averaging $\log r_t$ under the new policy, and which by averaging $-\log r_t$ under the old one? Write both in the
notation $p\kl{p}q$.
\end{exercise}

\appendix
\renewcommand{\thesection}{\Alph{section}}
\renewcommand{\thefigure}{\Alph{section}.\arabic{figure}}
\setcounter{figure}{0}

\section{Why ratios, and not differences?}\label{app:ratios}

The note I left myself was: \emph{what is the fundamental thing going on, and why are ratios more natural than standard
differences?} There are several answers, and they are one answer: $\log p-\log q$ is the right subtraction for
probabilities.

\begin{itemize}
\item \textbf{Evidence multiplies.} For $n$ independent samples the likelihood ratio of the whole dataset is the
product of the per-sample ratios, so its log is a sum, $S_n=\sum_{i\le n}\ell(x_i)$. KL is the expected evidence per
sample: $S_n$ drifts upward at slope $p\kl{p}q$ (Figure~\ref{fig:evidence}). Raw differences $p(x)-q(x)$ do not
compose across observations at all.
\item \textbf{Coordinates do not matter.} Under a change of variables $y=g(x)$ both densities pick up the same Jacobian;
in the ratio it cancels, in $p(x)-q(x)$ it does not. KL is a statement about distributions, not coordinates.
\item \textbf{Bayes is additive in log-ratios.} Posterior odds are prior odds times the likelihood ratio, so $\ell(x)$
is how far one observation moves your log-odds toward $p$, Good and Turing's \emph{weight of evidence}.
\item \textbf{Hypothesis testing.} By the Neyman--Pearson lemma, thresholding the likelihood ratio is the most powerful
test of $p$ against $q$. By Stein's lemma, with $n$ samples from $p$ the best achievable probability of wrongly
accepting $q$ decays like $e^{-n\DKL(p\Vert q)}$: KL is the exponent of distinguishability.
\item \textbf{Changing measure needs a ratio.} Importance sampling uses $\E_p[f(X)]=\E_q\big[\tfrac{p(X)}{q(X)}f(X)\big]$.
That is PPO's ratio, \textbf{new over old}, reweighting samples collected by the old policy. Averaged under the new
policy, $\log r$ gives $\DKL(\pi_\theta\Vert\pi_{\mathrm{old}})$; on the old policy's samples, $\E_{\mathrm{old}}[-\log r]$ gives
$\DKL(\pi_{\mathrm{old}}\Vert\pi_\theta)$, the direction TRPO constrains.
\end{itemize}

\begin{figure}[t]
\centering
\begin{tikzpicture}
\begin{axis}[width=11cm,height=5.4cm,xmin=0,xmax=200,ymin=-5,ymax=80,axis lines=left,
  xlabel={number of samples $n$ from $p$},ylabel={$S_n=\sum_{i\le n}\ell(x_i)$ (nats)},
  label style={font=\small},tick label style={font=\small},
  pathstyle/.style={dcolor!55,thick},
  legend style={draw=none,font=\small,at={(0.03,0.97)},anchor=north west}]
  \addplot[black,very thick,dashed,domain=0:200]{0.3181*x};
  \addlegendentry{$n\cdot\big(p\kl{p}q\big)=0.318\,n$}
  % seed 1
\addplot[pathstyle] coordinates {(0,0) (2,-0.02) (4,1.14) (6,2.08) (8,2.30) (10,3.67) (12,4.63) (14,6.01) (16,6.44) (18,5.65) (20,7.01) (22,7.81) (24,8.84) (26,9.81) (28,11.01) (30,11.88) (32,13.07) (34,14.00) (36,15.37) (38,16.14) (40,17.38) (42,17.29) (44,18.37) (46,18.83) (48,19.80) (50,20.35) (52,21.08) (54,21.06) (56,21.78) (58,22.96) (60,23.94) (62,24.73) (64,25.27) (66,26.23) (68,26.48) (70,27.50) (72,28.86) (74,29.33) (76,30.05) (78,31.29) (80,29.66) (82,30.47) (84,31.65) (86,30.76) (88,31.68) (90,32.47) (92,33.85) (94,33.95) (96,34.33) (98,35.16) (100,36.13) (102,35.99) (104,37.21) (106,38.10) (108,39.17) (110,39.82) (112,40.75) (114,41.94) (116,42.66) (118,42.85) (120,42.96) (122,42.97) (124,44.29) (126,45.66) (128,46.84) (130,47.49) (132,48.82) (134,49.64) (136,50.79) (138,51.72) (140,52.63) (142,53.65) (144,54.88) (146,54.54) (148,55.75) (150,55.86) (152,57.23) (154,57.67) (156,58.14) (158,58.94) (160,57.54) (162,58.38) (164,58.98) (166,59.72) (168,60.36) (170,61.48) (172,61.31) (174,61.23) (176,60.51) (178,61.49) (180,62.87) (182,64.23) (184,63.16) (186,64.41) (188,63.06) (190,63.91) (192,64.98) (194,65.53) (196,66.75) (198,67.31) (200,68.18)};
% seed 7
\addplot[pathstyle] coordinates {(0,0) (2,1.26) (4,2.59) (6,3.64) (8,4.49) (10,5.45) (12,6.77) (14,6.84) (16,8.04) (18,7.21) (20,8.22) (22,9.57) (24,10.70) (26,11.99) (28,12.11) (30,12.84) (32,13.88) (34,15.21) (36,16.43) (38,17.39) (40,18.12) (42,19.24) (44,19.72) (46,20.47) (48,20.30) (50,21.43) (52,22.72) (54,23.04) (56,23.93) (58,24.49) (60,25.23) (62,26.34) (64,27.05) (66,27.98) (68,27.19) (70,27.76) (72,28.24) (74,25.89) (76,27.03) (78,27.59) (80,28.51) (82,29.80) (84,30.09) (86,31.27) (88,31.11) (90,32.05) (92,31.83) (94,31.72) (96,32.70) (98,32.47) (100,33.73) (102,34.92) (104,35.81) (106,36.97) (108,37.95) (110,38.71) (112,39.21) (114,39.88) (116,41.22) (118,41.47) (120,41.66) (122,42.67) (124,43.30) (126,44.63) (128,45.89) (130,47.23) (132,48.49) (134,49.54) (136,49.37) (138,50.64) (140,51.70) (142,52.99) (144,50.65) (146,51.54) (148,52.84) (150,54.00) (152,55.25) (154,54.38) (156,55.64) (158,57.01) (160,55.52) (162,56.01) (164,57.05) (166,57.33) (168,57.58) (170,58.78) (172,57.02) (174,57.18) (176,57.55) (178,58.39) (180,59.76) (182,60.90) (184,61.40) (186,62.34) (188,60.41) (190,61.45) (192,62.65) (194,63.86) (196,63.52) (198,64.42) (200,64.60)};
% seed 2026
\addplot[pathstyle] coordinates {(0,0) (2,0.86) (4,0.77) (6,1.97) (8,2.75) (10,3.54) (12,3.83) (14,4.55) (16,5.22) (18,5.34) (20,5.46) (22,5.99) (24,7.16) (26,6.56) (28,1.97) (30,3.29) (32,4.07) (34,5.00) (36,5.47) (38,6.84) (40,7.88) (42,8.89) (44,9.76) (46,10.61) (48,11.31) (50,10.83) (52,11.81) (54,12.71) (56,14.08) (58,15.36) (60,15.83) (62,17.14) (64,18.05) (66,19.24) (68,19.88) (70,20.80) (72,18.91) (74,19.68) (76,20.79) (78,21.66) (80,21.42) (82,22.55) (84,23.91) (86,24.61) (88,25.45) (90,24.92) (92,25.81) (94,25.89) (96,26.66) (98,26.71) (100,26.23) (102,26.74) (104,27.58) (106,26.45) (108,26.45) (110,27.79) (112,28.66) (114,29.01) (116,29.95) (118,31.34) (120,30.70) (122,30.91) (124,31.56) (126,30.60) (128,31.23) (130,32.34) (132,32.62) (134,33.50) (136,34.78) (138,35.50) (140,35.32) (142,36.64) (144,35.97) (146,35.51) (148,36.50) (150,37.34) (152,38.71) (154,38.46) (156,39.52) (158,39.81) (160,40.66) (162,41.77) (164,43.06) (166,43.81) (168,45.08) (170,46.04) (172,46.74) (174,48.07) (176,49.29) (178,50.64) (180,50.66) (182,50.99) (184,51.45) (186,52.73) (188,52.96) (190,52.00) (192,51.99) (194,52.96) (196,54.28) (198,55.35) (200,56.22)};

  \addlegendentry{three runs}
\end{axis}
\end{tikzpicture}
\caption{Evidence adds. Cumulative log-likelihood ratio for $p=\mathcal N(0,1)$ against $q=\mathcal N(0,2^2)$ on samples
from $p$. Single steps can be negative (Figure~\ref{fig:pointwise}), but the sum drifts up at slope $p\kl{p}q\approx0.32$
nats per sample. Paths generated by \texttt{evidence\_paths.py}.}
\label{fig:evidence}
\end{figure}

Differences of probabilities do give real distances, such as total variation $\tfrac12\sum_x\lvert p(x)-q(x)\rvert$, and
Pinsker's inequality bounds it by KL: $\mathrm{TV}(p,q)\le\sqrt{\DKL(p\Vert q)/2}$. They just lack the properties above,
which is why ratios keep appearing everywhere in ML.

\setcounter{figure}{0}
\section{The classical bits view, side by side}

The textbook route starts from surprisal. Seeing $x$ with probability $p(x)$ costs $-\log_2 p(x)$ bits; for example
$p(x)=2^{-3}$ costs $3$ bits (Figure~\ref{fig:surprisal}). Averaging gives entropy. With the sampling distribution as a
subscript, as everywhere in this chapter,
\[
H_p(q):=\E_{X\sim p}[-\log q(X)]\quad\text{(cross-entropy)},\qquad H_p(p)=\E_{X\sim p}[-\log p(X)]\quad\text{(entropy)}.
\]
$H_p(q)$ is the expected code length when the data come from $p$ but the code was built for $q$; $H_p(p)$ is the best
possible. Their difference is the classical KL, the \textbf{excess bits} of using the wrong code:
\[
p\kl{p}q\;=\;H_p(q)-H_p(p)\;=\;\E_p[-\log q(X)]-\E_p[-\log p(X)]\;=\;\E_p[\log p(X)-\log q(X)].
\]

\begin{figure}[t]
\centering
\begin{tikzpicture}
\begin{axis}[width=8.6cm,height=5.2cm,xmin=0,xmax=1.05,ymin=0,ymax=6.2,axis lines=left,
  xlabel={$p(x)$},ylabel={$-\log_2 p(x)$ (bits)},label style={font=\small},tick label style={font=\small},
  xtick={0.125,0.25,0.5,1},xticklabels={$\frac18$,$\frac14$,$\frac12$,$1$},ytick={0,1,2,3,4,5,6},clip=false]
  \addplot[black,very thick,domain=0.0156:1,samples=200]{-ln(x)/ln(2)};
  \addplot[only marks,mark=*,mark size=1.8pt,dcolor] coordinates {(0.125,3) (0.5,1) (1,0)};
  \draw[dcolor,densely dotted] (axis cs:0.125,0) -- (axis cs:0.125,3) -- (axis cs:0,3);
  \draw[dcolor,densely dotted] (axis cs:0.5,0) -- (axis cs:0.5,1) -- (axis cs:0,1);
  \node[font=\small,text=dcolor,anchor=west] at (axis cs:0.15,3.25) {$p(x)=2^{-3}\Rightarrow3$ bits of surprise};
  \node[font=\small,text=dcolor,anchor=west] at (axis cs:0.52,1.3) {a fair coin flip: $1$ bit};
\end{axis}
\end{tikzpicture}
\caption{Surprisal, from page~5 of my notes: a bit of information is $-\log_2 p(x)$ when $X=x$. Rare outcomes cost more
bits; a sure outcome costs none.}
\label{fig:surprisal}
\end{figure}

The two readings are the same number thought of two ways.
\begin{itemize}
\item \textbf{Bits:} a difference of \emph{costs}. Bigger is worse, so every term carries a minus sign, and we subtract
the optimal cost from the suboptimal one.
\item \textbf{Log-likelihood:} a difference of \emph{plausibilities}. Bigger is better, there are no minus signs, and we
subtract the baseline model's log-likelihood from the true model's, on the true model's data.
\end{itemize}
The bits view is the right one when compression is the goal. For learning, where everything we optimize is already a
log-likelihood, the second is far easier to hold in one's head. They are the same subtraction.

\begin{thebibliography}{99}\small
\bibitem{akaike1973} H.~Akaike. Information theory and an extension of the maximum likelihood principle.
  \emph{2nd International Symposium on Information Theory}, 1973.
\bibitem{bishop2006} C.~M.~Bishop. \emph{Pattern Recognition and Machine Learning}. Springer, 2006.
\bibitem{bishop2024} C.~M.~Bishop and H.~Bishop. \emph{Deep Learning: Foundations and Concepts}. Springer, 2024.
\bibitem{csiszar1975} I.~Csisz\'ar. $I$-divergence geometry of probability distributions and minimization problems.
  \emph{Annals of Probability}, 3(1):146--158, 1975.
\bibitem{etz2018} A.~Etz. Some technical notes on Kullback--Leibler divergence. Blog post, 2018.
\bibitem{goodfellow2016} I.~Goodfellow, Y.~Bengio and A.~Courville. \emph{Deep Learning}. MIT Press, 2016.
\bibitem{hastie2009} T.~Hastie, R.~Tibshirani and J.~Friedman. \emph{The Elements of Statistical Learning}, 2nd ed.
  Springer, 2009.
\bibitem{hogg2013} R.~V.~Hogg, J.~W.~McKean and A.~T.~Craig. \emph{Introduction to Mathematical Statistics}, 7th ed.
  Pearson, 2013.
\bibitem{koller2009} D.~Koller and N.~Friedman. \emph{Probabilistic Graphical Models}. MIT Press, 2009.
\bibitem{kullback1959} S.~Kullback. \emph{Information Theory and Statistics}. Wiley, 1959.
\bibitem{kullback1951} S.~Kullback and R.~A.~Leibler. On information and sufficiency. \emph{Annals of Mathematical
  Statistics}, 22(1):79--86, 1951.
\bibitem{mackay2003} D.~J.~C.~MacKay. \emph{Information Theory, Inference, and Learning Algorithms}. Cambridge
  University Press, 2003.
\bibitem{ssbd2014} S.~Shalev-Shwartz and S.~Ben-David. \emph{Understanding Machine Learning}. Cambridge University
  Press, 2014.
\end{thebibliography}

\end{document}
